\chapter{Dateien}\label{kap:dateien}
\index{Datei}

\produkt{} speichert seine Arbeit in Dokumenten mit der Endung \datei{.rttf}. Daneben liest und
schreibt es STEP und schreibt STL. Dieses Kapitel beschreibt, was in einem Dokument steht, wie
Sie Dateien öffnen, speichern, importieren und exportieren und welche Beispielmodelle es gibt.
Den genauen Aufbau des Dateiformats beschreibt \autoref{kap:dateiformat}.

\begin{referenz}{>{\raggedright\arraybackslash}p{3.8cm}>{\centering\arraybackslash}p{1.2cm}>{\centering\arraybackslash}p{1.4cm}X}
\toprule
\textbf{Format} & \textbf{Lesen} & \textbf{Schreiben} & \textbf{Hinweise} \\
\midrule
\endhead
\datei{.rttf} (Raytatouille Fusion file) & \checkmark & \checkmark & das Dokument mit allem,
  was dazugehört (\autoref{sec:rttf}) \\
\datei{.nfcad} & \checkmark & (\checkmark) & dasselbe Format aus der Zeit vor Version 1.0;
  wird geöffnet, \gui{Save As...} schlägt \datei{.rttf} vor \\
STEP (\datei{.step}, \datei{.stp}) & \checkmark & \checkmark & exakte Geometrie; Öffnen
  ersetzt, Importieren ergänzt das Dokument; Export als AP214
  (\autoref{sec:step-import}, \autoref{sec:step-export}) \\
STL (\datei{.stl}) & -- & \checkmark & binär, aus den Dreiecksnetzen der Darstellung
  (\autoref{sec:stl-export}) \\
\datei{.nfplugin} & \checkmark & -- & Plugin-Paket; Öffnen installiert das Plugin
  (\autoref{kap:plugin-verwaltung}) \\
\bottomrule
\end{referenz}

Optische Systeme im Format von Raytatouille (\datei{.rtt.json}) und Glaskataloge
(\datei{.agf}) liest nicht das CAD-Programm, sondern das Plugin \plugin{Raytracer}
(\autoref{kap:raytracer}).

\section{Das Dokument (.rttf)}\label{sec:rttf}
\index{rttf}\index{Dokument}

Ein Dokument ist eine einzige Datei. Sie enthält alles, was nötig ist, um die Arbeit auf einem
anderen Rechner fortzusetzen -- auch die Daten, die Plugins für ihre Rechnungen brauchen. Die
Endung \datei{.rttf} steht für \enquote{Raytatouille Fusion file}.

\subsection{Was im Dokument steht}\label{sec:rttf-inhalt}

\begin{referenz}{>{\raggedright\arraybackslash}p{3.6cm}X}
\toprule
\textbf{Inhalt} & \textbf{Einzelheiten} \\
\midrule
\endhead
Geometrie & jedes Objekt \emph{exakt}: Grad, Knotenvektor, Kontrollpunkte und Gewichte der
  NURBS, bei getrimmten Flächen die Randkurven im Parameterraum, bei Hüllen die Liste ihrer
  Flächen. Keine Dreiecksnetze -- die Darstellung wird beim Öffnen neu berechnet. Die
  Topologie der Hüllen (welche Flächen an welcher Kante zusammenstoßen) wird beim Öffnen durch
  Zusammennähen wiederhergestellt. Einheit: Millimeter. \\
Objektdaten & Id, Name, Ebene, Sichtbarkeit, gegebenenfalls eigene Farbe, Gruppe und freie
  Attribute (\autoref{kap:organisieren}). \\
Ebenen & Name, Farbe, Sichtbarkeit (\autoref{sec:ebenen}). \\
Gruppen & Name, übergeordnete Gruppe, Sichtbarkeit -- also der ganze Modellbaum
  (\autoref{sec:gruppen}). \\
Plugin-Parameter & die Werte im Dock \gui{Plugin}, für jedes Plugin getrennt. Werte eines
  Plugins, das auf dem aktuellen Rechner nicht installiert ist, bleiben beim Speichern
  unverändert erhalten. \\
Workflow & der Knoten-Workflow mit Blöcken, Verbindungen, Zweigen, Gruppen und Varianten,
  einschließlich der Daten, die Blöcke aus Dateien übernommen haben
  (\autoref{sec:glaskataloge}). \\
Gespeicherte Ergebnisse & das neueste Ergebnis jedes Anzeigefensters des Workflows: Bilder,
  Kurven, Tabellen, Balken -- jeweils mit einem Fingerabdruck dessen, woraus es berechnet
  wurde (\autoref{sec:gespeicherte-ergebnisse}). \\
\bottomrule
\end{referenz}

\paragraph{Nicht im Dokument} stehen: der Zoom und die Blickrichtung der Ansichten, die
Darstellungsmodi, die Auswahl, die Liste der Rückgängig-Schritte, die Fensteranordnung (sie
gehört zu den Programmeinstellungen, \autoref{sec:einstellungen}) und die übrigen Ergebnisse
im Dock \gui{Results}, etwa Strahlen und Felder auf Flächen. Diese werden beim nächsten Rechnen
neu erzeugt.

\paragraph{Technisches.} Das Dokument ist eine UTF-8-Textdatei im JSON-Format. Sie wird
\emph{atomar} geschrieben: Das Programm schreibt zuerst eine temporäre Datei und ersetzt die
alte erst, wenn alles geschrieben ist. Bricht das Speichern ab (volle Festplatte, Absturz),
bleibt die vorige Fassung unversehrt. Das Format trägt eine Versionsnummer; ältere Fassungen
werden gelesen, eine neuere als die vom Programm unterstützte wird mit der Meldung
\enquote{document version \dots{} is newer than this program supports} abgelehnt.

\begin{hinweis}
Dokumente mit gespeicherten Bildern können groß werden, weil die Bildwerte mitgespeichert
werden (je Pixel vier Byte, in der Datei als Base64-Text). Bei sehr großen Kamerabildern
lohnt es sich, nicht mehr benötigte Anzeigeblöcke aus dem Workflow zu entfernen.
\end{hinweis}

\subsection{Ältere Dokumente (.nfcad)}\label{sec:nfcad}

Dateien mit der Endung \datei{.nfcad} stammen aus der Zeit vor Version~1.0, als das Programm
noch Nurbforge CAD hieß. Ihr Inhalt hat dasselbe Format; \produkt{} öffnet sie ohne
Umwandlung. \menu{File,Save} schreibt ein geöffnetes \datei{.nfcad}-Dokument unter seinem
bisherigen Namen zurück. \menu{File,Save As...} schlägt denselben Namen mit der Endung
\datei{.rttf} vor -- so stellen Sie ein Dokument auf die neue Endung um.

\section{Öffnen, Speichern, Neu}\label{sec:oeffnen-speichern}
\index{Öffnen}\index{Speichern}

\subsection{Öffnen}

\begin{schritte}
  \item \menu{File,Open...} (\keys{Strg+O}).
  \item Hat das aktuelle Dokument ungespeicherte Änderungen, fragt das Programm zuerst
    \enquote{The document has been modified. Save the changes?} -- \gui{Save} speichert,
    \gui{Discard} verwirft die Änderungen, \gui{Cancel} bricht das Öffnen ab.
  \item Der Dateidialog beginnt im Ordner des aktuellen Dokuments. Die Dateitypliste bietet
    \gui{All supported}, \gui{Raytatouille Fusion documents (*.rttf *.nfcad)} und \gui{STEP
    files}. Wählen Sie die Datei.
\end{schritte}

Das geöffnete Dokument ersetzt das bisherige vollständig; die Liste der Rückgängig-Schritte
und das Dock \gui{Results} werden geleert. Gespeicherte Ergebnisse kehren als Lauf
\enquote{Saved results} zurück, die Parameter der Plugins erscheinen in ihrem Dock, und ein
Workflow öffnet das Fenster \gui{Node workflow}. Sobald die Darstellung berechnet ist, zeigen
alle Ansichten das ganze Modell (\gui{Zoom extents}), und die Statusleiste meldet die Zahl der
Objekte und Dreiecke und die benötigte Zeit.

Lässt sich eine Datei nicht lesen, meldet ein Fenster \enquote{Cannot open \dots} mit dem
Grund, etwa \enquote{not a Raytatouille Fusion document} oder \enquote{invalid geometry:
\dots}. Das bisherige Dokument bleibt dann unverändert.

\begin{hinweis}
Eine Liste der zuletzt geöffneten Dateien bietet das Programm nicht. Ein Dokument öffnen Sie
über \menu{File,Open...}, unter Windows auch per Doppelklick im Explorer
(\autoref{sec:kommandozeile}).
\end{hinweis}

\subsection{Speichern}

\begin{itemize}
  \item \menu{File,Save} (\keys{Strg+S}) speichert unter dem bisherigen Namen. Hat das
    Dokument noch keinen (neues Dokument, geöffnete STEP-Datei), wirkt es wie \gui{Save As...}.
  \item \menu{File,Save As...} fragt nach einem Namen; gespeichert wird immer im Format
    \datei{.rttf}. Geben Sie keine Endung an, hängt das Programm \datei{.rttf} an.
\end{itemize}

Nach dem Speichern meldet die Statusleiste \enquote{Saved \textit{Dateiname}}, und das
Sternchen im Fenstertitel verschwindet. Scheitert das Speichern (Schreibschutz, voller
Datenträger), erscheint \enquote{Cannot save \dots} mit dem Grund.

\paragraph{Wann gilt ein Dokument als geändert?} Wenn es einen Bearbeitungsschritt gab
(\autoref{sec:undo}), wenn Plugin-Parameter geändert wurden, wenn sich der Workflow geändert
hat oder wenn eine Rechnung neue Ergebnisse für die Anzeigefenster des Workflows geliefert
hat. Nur dann fragt das Programm beim Schließen, beim Öffnen und bei \gui{New} nach dem
Speichern.

\begin{warnung}
Ein- und Ausblenden sowie Änderungen an Ebenen (neue Ebene, Name, Farbe) und Gruppennamen
gelten \emph{nicht} als Änderung. Sie werden zwar beim Speichern mitgeschrieben, lösen aber
keine Rückfrage aus. Haben Sie nur solche Dinge geändert, speichern Sie selbst.
\end{warnung}

\subsection{Neu}

\menu{File,New} (\keys{Strg+N}) beginnt ein leeres Dokument \enquote{Untitled} mit der Ebene
\gui{Default}, ohne Plugin-Parameter und ohne Workflow -- nach Rückfrage, falls das bisherige
ungespeicherte Änderungen hat.

\section{STEP importieren und öffnen}\label{sec:step-import}
\index{STEP}\index{Import}

STEP (ISO 10303) ist das übliche Austauschformat für exakte CAD-Geometrie. \produkt{} liest
STEP-Dateien der Anwendungsprotokolle AP203 und AP214 mit B-Spline- und analytischen Flächen
(Ebenen, Zylinder, Kegel, Kugeln, Tori \dots) und wandelt sie in NURBS-Hüllen um, in
Millimetern. Jeder Körper der Datei wird ein Objekt vom Typ \gui{Shell}; sein Name kommt aus
der STEP-Datei, oder er heißt nach der Datei mit einer laufenden Nummer.

Es gibt zwei Wege, eine STEP-Datei zu lesen:

\begin{referenz}{>{\raggedright\arraybackslash}p{3.8cm}X}
\toprule
\textbf{Weg} & \textbf{Wirkung} \\
\midrule
\endhead
\menu{File,Open...} (\keys{Strg+O}) & \emph{ersetzt} das Dokument durch den Inhalt der
  STEP-Datei (vorher Rückfrage bei Änderungen). Das neue Dokument hat noch keinen Namen;
  \gui{Save} fragt nach einem \datei{.rttf}-Namen, die STEP-Datei wird nicht überschrieben. \\
\menu{File,Import STEP...} (\keys{Strg+I}) & \emph{ergänzt} das aktuelle Dokument um den
  Inhalt der STEP-Datei. Das Einlesen ist ein Rückgängig-Schritt (\enquote{Import STEP}). \\
\bottomrule
\end{referenz}

In beiden Fällen bilden alle Körper der Datei eine Gruppe im Modellbaum, die den Dateinamen
ohne Endung trägt -- aus \datei{demo-parts.step} wird die Gruppe \enquote{demo-parts}.

\paragraph{Ablauf.} Das Lesen geschieht im Hintergrund; die Statusleiste zeigt
\enquote{Reading \textit{Datei} \dots}, danach \enquote{Tessellating \dots}, solange die
Darstellung berechnet wird, und schließlich die Zahl der Objekte und Dreiecke mit der
benötigten Zeit. Beim Öffnen und Importieren zeigt das Programm danach in allen Ansichten das
ganze Modell.

\paragraph{Warnungen und Fehler.} Elemente, die das Programm nicht umsetzen kann, werden
übersprungen. Ein Fenster nennt dann \enquote{\textit{n} warning(s) while reading:} mit den
ersten zwanzig Warnungen. Lässt sich die Datei gar nicht lesen, erscheint \enquote{Cannot read
\dots} mit dem Grund. Prüfen Sie nach Warnungen, ob alle Teile vollständig sind -- etwa im
Modellbaum und an den Eigenschaften der Hüllen (\gui{Closed} bei Volumenkörpern).

\begin{tipp}
Um mehrere Baugruppen in ein Dokument zu holen, öffnen Sie die erste mit \gui{Open...} und
holen Sie die weiteren mit \gui{Import STEP...} dazu. Jede Datei wird eine eigene Gruppe, die
Sie im Modellbaum einzeln ein- und ausblenden können.
\end{tipp}

\section{STEP exportieren}\label{sec:step-export}
\index{STEP!Export}\index{Export}

\menu{File,Export STEP...} schreibt Geometrie als STEP-Datei nach AP214, in Millimetern.

\begin{schritte}
  \item Wählen Sie die Objekte aus, die exportiert werden sollen. Ohne Auswahl exportiert das
    Programm alle \emph{sichtbaren} Objekte; ausgeblendete und solche auf ausgeblendeten
    Ebenen oder in ausgeblendeten Gruppen bleiben dann außen vor.
  \item \menu{File,Export STEP...}, Dateiname eingeben. Ohne Endung hängt das Programm
    \datei{.step} an.
  \item Die Statusleiste meldet \enquote{Exported \textit{n} object(s) to \textit{Datei}}.
\end{schritte}

Exportiert werden Hüllen (Volumenkörper und offene Flächenverbände), getrimmte und untrimmte
Flächen -- jede Fläche als eigene Hülle --, und zwar exakt, ohne Annäherung. Jedes Objekt
behält seinen Namen (ohne Namen \enquote{shell\_\textit{Id}}). \emph{Nicht} exportiert werden
Punkte und Kurven sowie Ebenen, Gruppen, Farben und Attribute. Enthält die Auswahl nichts
Exportierbares, meldet die Statusleiste \enquote{Nothing to export (STEP: solids and surfaces,
STL: objects with faces)}.

\begin{hinweis}
Teilobjekte (einzelne Flächen einer Hülle) werden nicht einzeln exportiert. Um nur eine Fläche
einer Hülle weiterzugeben, zerlegen Sie die Hülle zuerst mit \cmd{Explode}
(\autoref{kap:modellieren}).
\end{hinweis}

\section{STL exportieren}\label{sec:stl-export}
\index{STL}

\menu{File,Export STL...} schreibt die Dreiecksnetze der ausgewählten (ohne Auswahl: aller
sichtbaren) Objekte mit Flächen in \emph{eine} binäre STL-Datei, in Millimetern. Ohne Endung
hängt das Programm \datei{.stl} an. Punkte und Kurven haben keine Dreiecke und fehlen.

Die Dreiecke sind die der Bildschirmdarstellung: Das STL ist so genau wie die Darstellung in
den Ansichten, nicht wie die exakte Geometrie. Für 3D-Druck, Visualisierung und
Netzwerkzeuge ist das meist ausreichend; für die Fertigung oder für andere CAD-Programme
verwenden Sie STEP.

\section{Öffnen per Doppelklick und Kommandozeile}\label{sec:kommandozeile}
\index{Kommandozeile}\index{Doppelklick}

\paragraph{Doppelklick.} Unter Windows verknüpft das Installationsprogramm die Endungen
\datei{.rttf} und \datei{.nfcad} mit \produkt{}; ein Doppelklick im Explorer startet das
Programm mit dem Dokument. Unter Linux meldet der Anwendungseintrag das Programm für
STEP-Dateien an, sodass es im Dateimanager unter \enquote{Öffnen mit} erscheint. Unter macOS
ist das Programm mit keiner Endung verknüpft; öffnen Sie Dokumente dort über
\menu{File,Open...} oder aus dem Terminal.

\paragraph{Kommandozeile.} Das Programm nimmt eine Datei als Argument:

\begin{lstlisting}
RaytatouilleFusion [Optionen] [Datei]
\end{lstlisting}

Der Programmaufruf lautet je nach System:
\begin{lstlisting}
macOS:   /Applications/RaytatouilleFusion.app/Contents/MacOS/RaytatouilleFusion
Windows: "C:\Program Files\Raytatouille Fusion\bin\RaytatouilleFusion.exe"
Linux:   RaytatouilleFusion
\end{lstlisting}

\begin{referenz}{>{\raggedright\arraybackslash}p{4.4cm}X}
\toprule
\textbf{Argument} & \textbf{Wirkung} \\
\midrule
\endhead
\textit{Datei} & ein Dokument (\datei{.rttf}, \datei{.nfcad}) oder eine STEP-Datei, die geöffnet
  wird, oder ein Plugin-Paket (\datei{.nfplugin}), das installiert wird. \\
\code{-{}-display-mode} \textit{Modus} & Darstellungsmodus aller Ansichten beim Start:
  \code{wireframe}, \code{surfaces} oder \code{shaded}. \\
\code{-{}-benchmark} \textit{Bilder} & Messlauf: öffnet die Datei, dreht die perspektivische
  Ansicht über die angegebene Zahl von Bildern, gibt die Zeiten aus und beendet das Programm. \\
\code{-{}-help}, \code{-h} & listet die Optionen auf. \\
\code{-{}-version}, \code{-v} & gibt die Versionsnummer aus. \\
\bottomrule
\end{referenz}

Beispiel (macOS):
\begin{lstlisting}
/Applications/RaytatouilleFusion.app/Contents/MacOS/RaytatouilleFusion --display-mode wireframe ~/Modelle/objektiv.rttf
\end{lstlisting}

Lässt sich ein beim Start angegebenes Dokument nicht öffnen, zeigt das Programm die
Fehlermeldung und beendet sich danach.

\section{Beispielmodelle}\label{sec:beispiele}
\index{Beispiele}

Zum Ausprobieren gibt es Beispiele. Die Beispiele im Ordner \datei{examples} gehören nicht zur
Installation; Sie finden sie im Quelltext-Archiv des Projekts auf GitHub
(\url{https://github.com/PhotonPower/Raytatouille-Fusion}, Ordner \datei{examples}).

\begin{referenz}{>{\raggedright\arraybackslash}p{5.6cm}X}
\toprule
\textbf{Datei} & \textbf{Inhalt} \\
\midrule
\endhead
\datei{examples/demo-parts.step} & zwölf Volumenkörper -- Platten mit Löchern (boolesche
  Differenzen), Kugeln, Tori, Kegel, Zylinder -- in einem Raster. Gut zum Üben von Navigation,
  Auswahl, Teilobjekten und Darstellungsmodi (\autoref{fig:hauptfenster}). \\
\datei{examples/pathtracer/\allowbreak cornell-box.rttf} & Cornell-Box mit roter und grüner Wand, Lampe
  in der Decke, einem Block und einer Glaskugel; bereit für \gui{Camera image} des Path
  Tracers (\autoref{kap:pathtracer}). \\
\datei{examples/pathtracer/\allowbreak caustics.rttf} & Kaustiken: eine stark dispergierende Glaskugel,
  ein blauer Glasring und ein Band aus poliertem Gold unter einem Spot. \\
\datei{examples/pathtracer/\allowbreak prism.rttf} & ein Prisma aus Schwerflint fächert weißes Licht in
  31 Wellenlängen auf. \\
\datei{examples/pathtracer/\allowbreak depth-of-field.rttf} & fünf Kugeln aus verschiedenen Materialien,
  aufgenommen mit Schärfentiefe. \\
\datei{singlet\_optim.rtt.json} (im Raytracer-Plugin, Ordner \datei{examples}) & ein Raytatouille-System: eine
  Einzellinse zum Optimieren, mit konstanter Brechzahl (kein Glaskatalog nötig). Wird nicht
  mit \gui{Open...} geöffnet, sondern im Plugin \plugin{Raytracer} über \gui{System file} und
  \gui{Load system} (\autoref{kap:raytracer}). \\
\bottomrule
\end{referenz}

Das Raytracer-Beispiel wird mit dem Programm installiert, im Plugin-Ordner:
\begin{lstlisting}
macOS:   /Applications/RaytatouilleFusion.app/Contents/Resources/python/plugins/raytracer/examples
Windows: C:\Program Files\Raytatouille Fusion\python\plugins\raytracer\examples
Linux:   /usr/share/raytatouille-fusion/python/plugins/raytracer/examples
\end{lstlisting}

Weitere optische Systeme liegen dem Rechenkern Raytatouille bei; Systeme mit realen Gläsern
brauchen dessen Glaskataloge (\datei{.agf}) im Raytracer-Panel unter \gui{Glass catalogues}.

\section{Eingebettete Daten: Glaskataloge und andere Dateien}\label{sec:glaskataloge}
\index{Glaskatalog}

Ein Dokument soll auf jedem Rechner rechnen, auch dort, wo die Dateien fehlen, aus denen es
entstanden ist. Deshalb übernimmt der Workflow Daten aus Dateien in das Dokument:

\begin{itemize}
  \item \textbf{Glaskataloge.} Die Glaskataloge eines optischen Systems werden an einer Stelle
    verwaltet: in der Liste \gui{Glass catalogues} im Dock \gui{Plugin} beim
    \plugin{Raytracer}. Dazu kommen alle \datei{.agf}-Dateien neben der Systemdatei und in deren
    Unterordner \datei{catalogs}. Beim nächsten \gui{Check workflow} oder \gui{Run workflow}
    speichert der Workflow diese Kataloge im Dokument (in der Quelle des Systems). Danach
    rechnet die \datei{.rttf}-Datei auch ohne die \datei{.agf}-Dateien -- im Workflow ebenso
    wie im Raytracer-Panel, das die gespeicherten Kataloge benutzt, wenn eine Datei fehlt.
  \item \textbf{Dateien von Blöcken.} Dateien, die Blöcke des Workflows nennen -- die
    Konfiguration oder Jones-Tabelle eines Gitters, die Konfiguration eines Schichtstapels --,
    speichert \gui{Check workflow} (und \gui{Run workflow}) ebenfalls in den Blöcken und löscht
    die Pfade. Eine genannte Datei, die sich nicht lesen lässt, ergibt eine Warnung.
  \item \textbf{Ausnahme:} STEP-Dateien von CAD-Schichten in einer Konfiguration von
    FourierGrating
    bleiben außerhalb des Dokuments; der Check meldet sie als externe Dateien
    (\autoref{kap:fouriergrating}).
\end{itemize}

Fehlen beim Rechnen Gläser, nennen das Protokoll und die Statuszeile den Katalog, die
fehlenden Gläser und die Abhilfe. Einzelheiten in \autoref{kap:raytracer} und
\autoref{kap:workflow}.

\begin{tipp}
Bevor Sie ein Dokument weitergeben, führen Sie einmal \gui{Check workflow} aus und speichern
danach. Dann stecken Glaskataloge und Blockdateien im Dokument, und der Empfänger braucht nur
die \datei{.rttf}-Datei.
\end{tipp}

\section{Gespeicherte Ergebnisse}\label{sec:gespeicherte-ergebnisse}
\index{Ergebnis!gespeichert}\index{Fingerabdruck}

Rechnungen im Knoten-Workflow können Minuten bis Stunden dauern. Damit ihre Ergebnisse nicht
verloren gehen, speichert das Programm mit dem Dokument das \emph{neueste Ergebnis jedes
Anzeigefensters} des Workflows: Kamerabilder, Pupillenbilder, Kurven von Parameterstudien,
Tabellen und Balkendiagramme.

\paragraph{Fingerabdruck.} Zu jedem gespeicherten Ergebnis gehört ein Fingerabdruck dessen,
woraus es berechnet wurde -- Workflow und Einstellungen. Nach dem Öffnen stehen die
Ergebnisse im Dock \gui{Results} als Lauf \enquote{Saved results} (Typ \enquote{document}),
und ihre Anzeigefenster zeigen sie wieder an. Die Workflow-Maschine vergleicht den
Fingerabdruck mit dem, was sie jetzt berechnen würde:

\begin{itemize}
  \item Sind Workflow und Einstellungen unverändert, gilt das Ergebnis als aktuell: Der Block
    trägt ein grünes Häkchen, und \gui{Run workflow} rechnet ihn nicht noch einmal (in der
    Statuszeile \enquote{(unchanged)}).
  \item Ändern Sie etwas, auf dem das Ergebnis beruht, wird es beim nächsten Lauf neu
    berechnet.
\end{itemize}

\paragraph{Was nicht gespeichert wird.} Zwischenergebnisse -- Strahlenbündel, Felder --
werden nicht im Dokument gespeichert, ebenso wenig die Ergebnisse von Plugin-Aktionen außerhalb
der Anzeigefenster des Workflows (z.\,B.\ \gui{Trace rays} des Raytracers). Bildwerte werden
mit einfacher Genauigkeit (32 Bit) gespeichert, Kurven und Balken mit doppelter.

\paragraph{Geändert durch Rechnen.} Liefert ein Lauf neue Ergebnisse für die Anzeigefenster,
gilt das Dokument als geändert, auch wenn Sie sonst nichts angefasst haben -- das Programm
fragt beim Schließen also nach dem Speichern, damit die Rechenzeit nicht verloren geht.
